\chapter*{The asterisk}
\phantomsection
\addcontentsline{toc}{chapter}{The asterisk}
\markboth{The asterisk}{The asterisk}

I'm often jealous, and slightly in awe, of folks with strong and
detailed biographical memories. I can recall, for example, the feeling
of watching an Indonesian harbour appear out of the fog: small rocky
islands jutting out of a perfectly calm dark sea, not long after the sun
had risen. I was working on an 18-metre yacht out of Darwin, Australia,
but I have no idea which island it is that teases my memory, or when this
was.

I don't remember where we were coming from or where we were headed. I could
probably work out the year, but most other details may be irretrievable:
the captain's name or that of the yacht, though I'm strangely sure of
its size. I remember a married couple was also among the crew, and
perhaps one of my crew mates may have been named Chris. Or perhaps
Steve.

The plan, as I remember it, had sounded almost absurdly generous: a
leisurely pinball from island to island, starting in Ambon and heading
south-west, then up past Malaysia and Thailand, around India, through the
Suez, lingering in the Mediterranean, and eventually on to North America.
But the captain, an Australian who had never been outside Australia,
wanted to abandon the island-by-island route, get to Bali quickly, and
then hurry on to what he called civilization.

We argued, I know, about that change, and I was left somewhere on or
near Komodo Island. I recall somehow explaining to a man in a small
village that I wanted to leave the island. A wooden fishing boat took me
out into the strait. Eventually a large ferry appeared and took me on
board. And where did I go next? Bali? Flores? I certainly spent time on
each island. But the sequence is lost to me.

To tell the story at all, I have to reconstruct it. I have a few fixed
points, a foggy harbour, a quarrel, a boat, a ferry, and
a set of gaps between them. Some parts were witnessed. Some are inferred.
Some are just blank.

For one of those states, the inferred but unseen, historical linguists
use the asterisk \mentionhead{*}. If you have any association at all with
that mark, it's likely connected to footnotes or perhaps the marking of
birth dates. It's a symbol of additional information, a small detour from
the main text. But within the world of linguistics, this little star
became a mark for something missing.

August Schleicher became a linguist, or philologist, before dictionaries
traced the history of our words with the precision they have today,
before scholars confidently linked the English word \mention{father}
with German \mention{Vater} and Sanskrit \mention{pitar}, for instance.
The background to Schleicher's use of the asterisk begins with Sir
William Jones.

Jones was a British judge stationed in India in the late 18th century.
While studying Sanskrit, the ancient sacred language of the subcontinent,
he began to see a system in the parallels between Sanskrit, Greek, and
Latin, in words for family
members, body parts, and basic concepts that seemed to echo across
millennia and continents. He theorized that these languages descended
from a common ancestor, a prehistoric root from which the diverse
branches of a language family had grown.

Schleicher inherited Jones's way of seeing. He wanted a system that could
explain regularities in language change, and he wanted to show the
missing links in the chain. That need gave the asterisk its first
linguistic job: to mark a form that the scholar inferred, but hadn't
directly seen.

Take the seemingly unremarkable English word \mention{tooth}. The
inherited word belongs to a well-established Indo-European cognate set:
Old English \mention{tōþ}, from Proto-Germanic *\mention{tanþ-}
\ensuremath{\sim} *\mention{tunþ-}, beside Gothic \mention{tunþus}, is
cognate with Latin \mention{dēns}, genitive \mention{dentis}, Greek
\mention{odoús} / \mention{odṓn}, genitive \mention{odóntos}, Sanskrit
\mention{dánta-}, Lithuanian \mention{dantìs}, Welsh \mention{dant},
and Old Irish \mention{dét} \citep{OEDtooth2026}.

The details matter less here than the pattern. Several languages
preserve related forms, and the missing ancestor has to be inferred from
the family resemblance.

This is the kind of pattern Schleicher was after. While there's written
evidence for words like \mention{Zahn}, \mention{tunþus}, and
\mention{dánta-}, he needed a way to mark the hypothetical: a
reconstructed stem, *\mention{dant-}, from which they seemed to have
descended \citep[53]{Schleicher1874Compendium}. He wasn't afraid to
guess at its existence, but he wanted to be clear that it was only a
guess.

The asterisk made the guess visible. It told the reader: this form is
part of the argument, but the evidence for it is indirect. The symbol
stood at the point where observation gave way to inference.

That's close to what memory does, at least when memory is honest. I
have seen that Indonesian harbour. I remember the dark water and the
small rocky islands. The rest of the trip exists for me as a sequence of
reconstructed positions, some firmer than others. In Schleicher's work,
the asterisk held open that same distinction between what was witnessed
and what was inferred.

Within linguistics, though, the asterisk now usually does another job.
When I tell folks this, they imagine the linguist as a teacher marking
up student papers with asterisks. But that's not it at all. Often the
ungrammatical sentences linguists are \enquote{starring} are those
they've made up themselves. Rarely is it to call attention to the kinds
of mistakes that people make. In fact, it's often to illustrate the kinds
of mistakes people \emph{don't} make.

If this sounds bizarre, it is. The starred sentences are regularly so
wonderfully off-kilter that I sometimes marvel at the ingenuity needed
to come up with them. To give you a sense of what I mean,
(\ref{ex:weird-ungrammaticals-b}) presents six typical oddballs from
\textit{The Cambridge Grammar of the English Language}.

\ea \label{ex:weird-ungrammaticals-b}
    \ea[*]{\mention{He offered her the steak fiendishly hungry.}}
    \ex[*]{\mention{An apparently bird hit the car.}}
    \ex[*]{\mention{The news his guilt was devastating.}}
    \ex[*]{\mention{It had been for you.}}
    \ex[*]{\mention{She lives here any longer.}}
    \ex[*]{\mention{She wrote more plays than he wrote five novels.}}
    \z
\z

These aren't ordinary schoolroom errors. They're probes, constructed to
test where English allows a phrase to go, what kind of word can appear
where, or which pieces of meaning can be combined.

% EDITORIAL-HPC: This is the first place to plant the detector/category
% distinction. The asterisk turns a feeling into a claim, but it doesn't yet
% tell us whether the claim is about grammaticality proper, acceptability,
% processing, register, pedagogy, or some neighbouring coupling.
This book follows that feeling you get when you read such examples: the
little jolt that comes from reading a sentence English seems not to
allow. But the asterisk does more than report that feeling. It turns it
into a claim. So the small historical question opens into a larger one:
how did a mark for reconstructed words become a mark for impossible
sentences, and what kind of evidence can such a mark be taken to
represent?

The answer begins with Henry Sweet (1845--1912). While Schleicher is the
father of the asterisk in linguistics, he never used it to mean
`ungrammatical'. That use appears to be Sweet's innovation.

Sweet was a scholar of phonetics, Old English, grammar, and language
teaching. Keenly observant of the mismatch between traditional grammar
rules and the realities of everyday speech, he developed a descriptive
approach to grammar. His \textit{A new English grammar: Logical and
historical} moved away from didactic notions of right and wrong
usage. In the preface to that book, he wrote,

\begin{quote}
    As my exposition claims to be scientific, I confine myself
to the statement and explanation of facts, without attempting
to settle the relative correctness of divergent usages. If an
`ungrammatical' expression such as \mention{it is me} is in general
use among educated people, I accept it as such, simply
adding that it is avoided in the literary language. \phantom{~ }\hfill(p. xi)
\end{quote}

With this focus on the language in use, Sweet repurposed Schleicher's
symbol. Where Schleicher had marked words not directly observed, Sweet
used the asterisk to signal forms that could be constructed from English
words but not accepted as ordinary English sentences. An early clear use
appears when Sweet notes that from \mention{these tall men} we can infer
\mention{these men are tall}, but we can't make \mention{some Englishmen}
into *\mention{Englishmen are some}, or \mention{half the island} into
*\mention{the island was half} \citep[14]{Sweet}.

That's no small shift. A reconstructed word is missing because
history failed to preserve it. A starred sentence is missing for another
reason: an absence in the practice of speakers. It belongs to the
grammar as a boundary case, a shape that helps us see where the system
won't go.
% EDITORIAL-HPC: Watch "absence in the practice of speakers." The HPC book
% distinguishes low frequency, conditioned rarity, unattestedness, and failure
% of licensing. This paragraph should eventually avoid implying that not being
% observed is the same kind of absence as being ungrammatical.

Today, the asterisk is as likely to mark an ungrammatical sentence as a
reconstructed word. I used the symbol for years before I knew this
history, proof of how deeply it's woven into the field.

I didn't arrive at the problem through Schleicher or Sweet. I arrived
at it through teaching, and more particularly through the strange
institutional confidence with which some rules are treated as grammatical
boundaries. It feels like my own interest in grammar should have a single
beginning, but the details are again frustratingly elusive. I have
memories of hours, whole afternoons, spent on the sixth floor of the
Kinokuniya book store in Shinjuku. This was the English-language floor,
where you'd find the English novels along with translations of Japanese
novels.

There was a section for language-teaching materials, for learning
Japanese and for teaching English, and I made extensive use of both. I
remember being surprised at the number of people standing and reading
manga, magazines, and books in book shops and convenience stores.
Somebody told me it was so normal they had a noun for it:
\mention{tachiyomi}, reading while standing in the shop.

It wasn't clear to me what the social boundaries of \mention{tachiyomi}
were, but it seemed to license my going into Kinokuniya, picking up any
book in the Cambridge Applied Linguistics series, and reading until my
legs and back hurt from standing so long. I didn't yet think of that as
a grammatical question, but it was already a question about rules,
permission, tacit knowledge, and the difference between what people say
is allowed and what people actually allow.

There was no way I was going to be able to buy any of those tomes, which
typically sold for something like ¥8,000, a painful fraction of a
month's rent. Standing time was the only rent I could afford, and a few
chapters on language were what I got for it.

But none of this had been planned. I hadn't gone to Japan because I had a plan to teach English. I'd run out of money around the time I reached Singapore and was trying to get myself home. While I planned, an Australian backpacker sharing my room pointed out that Japan needed English teachers. I needed money, and I spoke English. We were compatible.

I arrived in Tokyo in late fall, 1993. I had no relevant training and little experience except for a few private English lessons I'd taught when I was in Chiang Mai, perhaps a year earlier. I had also arrived in Japan knowing literally nothing about its language, and I soon decided I was going to learn it faster and better than the other foreigners around me, a task at which I was only partially successful.

While I focused on memorizing vocabulary and sentence patterns, the
sheer difference between Japanese grammar and the grammar I knew
constantly nagged at me. Phrases were built differently, verb
conjugations seemed baffling, and particles added layers of meaning I
struggled to grasp.

Japanese didn't just give me new words for familiar sentences. It made
different structural possibilities visible, and made some familiar
English ways of packaging meaning feel less inevitable than they had
before. In those days in Japan, I treated that as a problem of learning
and teaching. It would take something else to turn my attention more to
grammar itself.

I moved back to Canada with my wife and toddler in 2003, just before the
first global SARS outbreak. I was looking for language teaching work in
Toronto, but students were fleeing and universities were laying off
English-language teachers. After four months of searching, I was finally
interviewed for a full-time position at Humber College, a position I
still hold today.

I remember, not long after arriving, perhaps a few months later, perhaps
a few years, a rising sense of confusion and embarrassment over FANBOYS,
a mnemonic for the \enquote{co-ordinating conjunctions} in English:
\mention{for}, \mention{and}, \mention{nor}, \mention{but},
\mention{or}, \mention{yet}, and \mention{so}. I was confused because I
didn't understand why this idea seemed so oddly central in the Humber
English-teaching ecosystem. I was embarrassed because I'd never
encountered it before.

I recall turning to \textit{The Cambridge Grammar of the English
Language} to try to sort it out. I'd seen a copy years earlier at the
Temple University library in Japan, but hadn't been able to make much of
it. Humber also had a copy, so I thought I'd give it another try. This
time, I was able to make more sense of it and realized that I needn't
have been embarrassed about FANBOYS.

The experience led me to start my blog \textit{English, Jack}, where the
first of roughly 800 posts, on Friday, July 28, 2006, was entitled
\enquote{The myth of FANBOYS}. It also led me into a deep dive of
\textit{The Cambridge Grammar}. It was around this time that I also
encountered Geoff Pullum's work on the philosophy of grammar, which
pushed me toward the foundational questions of what grammar is, how we
know it, and what kind of evidence a judgment provides.

FANBOYS mattered because it felt like an asterisk without the star. Here
was a rule treated as grammar, taught with institutional confidence, and
used to sort good sentences from bad ones. But the more I looked at it,
the less it seemed to name a real boundary in English. The question was
no longer whether a particular sentence sounded right. The question was
what kind of thing a grammatical boundary could be.
% EDITORIAL-HPC: Good place to distinguish usage-level grammatical boundaries
% from metalinguistic/pedagogical kinds. FANBOYS may be real as an institutional
% teaching object while failing as a grammatical category.

Over the next decade, I submerged myself in the study of English grammar,
figuring out how it works and writing about what I learned. But a larger
question started to press on me: what kind of thing is
(un)grammaticality? Is an ungrammatical sentence like a social faux pas,
a violation of an unwritten etiquette? Or is it more basic, a strain on
the machinery by which we make language intelligible?

That question is what the asterisk hides. In 1973, Fred Householder
\citep{Householder1973} pointed out that the asterisk makes
ungrammaticality look simpler than it is. Different speakers, or even
the same speaker in different situations, might have different judgments
on some sentences. The asterisk presents grammaticality as a simple yes
or no, when the thing it marks is more complex.

One small star can cover several claims. It can mean that no competent
speaker would produce the sentence. It can mean that some speakers reject
it and others accept it. It can mean that the example needs a special
context, or that the sentence is hard to process, or that it violates a
classroom rule, or that it sounds wrong in this register but ordinary in
another. The mark is slight. The work it does isn't.

% EDITORIAL-HPC: This is exactly right as a pre-theory description of asterisk
% practice, but ch. 13 must resolve it. Ordinary uses of "ungrammatical" are a
% family; grammaticality proper should not become a bag of unrelated wrongnesses.
Every asterisk in this book asks a version of that question. It may mark
a sentence nobody says, a sentence some people say, a sentence someone
has been taught to dislike, or a sentence that strains the machinery of
understanding. The star looks like a single judgment. The book begins
with the suspicion that it marks a family of them.
